

\chapter{Getting  started 1D}
\label{chap:Gettingstarted1D}
In this  section we  will  see,  step by  step,  how  to use  \TiberCAD{} to simulate numerically a  semiconductor  device.
As a very  simple example  we  will  refer  to  the  \textbf{Tutorial 0} (Si bulk) that you  can  find  in  the  \textit{Tutorials}  directory.



\section*{Step 1: Modeling the device}

As a first step, we have to model the device. To do so, you can use DEVISE module of ISE-TCAD 9.5 software package or GMSH program.
Here we'll see in details the procedure for GMSH.
There are two possible ways to use GMSH:

    
\begin{enumerate}
	\item Interactive, using the graphical interface
   \item Using a script file.
\end{enumerate}

In the following we'll see how to write a basic GMSH script (bulk.geo); for any details please refer to GMSH manual GMSH (http://geuz.org/gmsh/).

 

 

In a GMSH script, several variables can be defined and given a value in this way:

\begin{verbatim}
L = 1;
d = 0.01;
\end{verbatim}

these are valid GMSH variables: \textit{L} is just the length of the Si sample; \textit{d} is the value of a \textbf{characteristic}
\textbf{mesh length} (see below).
\begin{itemize}
	\item Definition of geometrical entities \textbf{Points}:
\end{itemize}


\begin{verbatim}
Point(1) = {0, 0, 0, d};
Point(2) = {L, 0, 0, d};
\end{verbatim}

In the definition of a geometrical point, the three first expressions inside the braces on the right hand side give the three X, Y and Z coordinates of the point ; the last expression \textit{(d)} sets the \textbf{characteristic mesh length} at that point, that is the \textbf{size} of a mesh element, defined as the length of the segment for a line segment, the radius of the circumscribed circle for a triangle and the radius of the circumscribed sphere for a tetrahedron.

Thus, the smaller is the value of \textit{d,} the greater is the mesh density close to that point. The size of the mesh elements will then be computed in GMSH by linearly interpolating these characteristic lengths in the whole mesh.

\begin{itemize}
	\item Definition of geometrical entity \textbf{Line}:
\end{itemize}
  
\begin{verbatim}
Line(1) = {1, 2};
\end{verbatim}
The two expressions inside the braces on the right hand side  give the identification numbers of the start and end points of the line.

\begin{itemize}
	\item Definition of the physical entity \textbf{Physical Line 1}:
\end{itemize}

\begin{verbatim}
Physical Line(1) = {1};
\end{verbatim}

The expression(s) inside the braces on the right hand side  give the identification numbers of all the geometrical lines that need to be grouped inside the \textbf{physical line}.
In this way, in general, \textbf{physical regions} are created which associate together geometrical regions, and then the related mesh elements, which share some common physical properties. It's only these physical regions which can be referred to outside GMSH. In \TiberCAD{}, this is done by associating one or more physical regions to a \textbf{TiberCAD{} region }through the keyword \textbf{mesh\_regions} (see in the following).

\begin{itemize}
\item Definition of two physical entities\textbf{ Physical Point}:  
\end{itemize}

\begin{verbatim}
Physical Point(1) = {1};
Physical Point(2) = {2}; 
\end{verbatim}

Beginning  from  GMSH v.2, it  is  possible, alternatively,   to  assign  more  convenient \textbf{Physical Names}  to  the  Physical entities, instead of  numerical IDs.  \textbf{Physical Names} consist of  strings  enclosed between  quotation marks.
The  syntax is the  following:

\begin{verbatim}
Physical Line("bulk") = {1}
............
Physical Point("Anode") = {1};
\end{verbatim}



    \textbf{N.B}.: In general, in a \textit{nD} simulation, \textit{(n-1)D }physical regions (points in 1D, lines in 2D, surfaces in 3D) are used by \TiberCAD{} to impose the required boundary conditions.
      Each  \textit{(n-1)D } physical region defined in this way in GMSH will be associated in \TiberCAD{} to a boundary condition region, through the keyword \textbf{BC\_reg\_numb}. Thus, in this case, Physical points \textbf{1} and \textbf{2} will be associated respectively to two BC regions (see in the following).






\section*{Step 2: Meshing the device}

The \textit{.geo} script file with the geometrical description can be run in GMSH, to display the modelled device and to mesh it through the GMSH graphical interface.
Alternatively, a \textit{non-interactive} mode is also available in GMSH, without graphical user interface. For example, to mesh this 1D tutorial in non-interactive mode, just type:

\textit{gmsh bulk.geo  -1 -o bulk.msh }

where \textit{bulk.geo}  is the geometrical description of the device with GMSH syntax;
\textit{-1} means 1D mesh generation;

some command line options are:

\textit{-1, -2, -3} to perform 1D, 2D or 3D mesh generation,

\textit{-o} mesh\_file\textit{.msh} to specify the name of the mesh file to be generated

In this way, a \textit{.msh} has been generated and is ready to be read in \TiberCAD{}. 





\section*{Step 3: TiberCAD{}  Input file  }

Now we have to write down the \TiberCAD{} \textbf{input file} (see \textit{bulk.tib} in  the  Tutorials).


\subsection* {1 - Definition of Device Regions}
First, we have to list all the \TiberCAD{} Regions present in our Device: a \TiberCAD{} \textbf{Region} is usually a section of the device featuring the same material and possibly the same doping.
\begin{verbatim}
$Device
{
Region bulk
{
mesh_regions = 1
material = Si
doping = 1e16 doping_type = donor
}

}
\end{verbatim}
The \TiberCAD{} \textbf{Region} bulk is made of Silicon and n-doped with a concentration $10^{16} cm^{-3}$ .

Through the keyword  \textbf{mesh\_regions} , one or more of the \textit{physical regions} (\textit{Physical Lines} in 1D, \textit{Physical Surfaces} in 2D, \textit{Physical Volumes} in 3D) previously defined in the GMSH mesh can be associated to the present  \TiberCAD{} \textbf{Region}.

With \textit{mesh\_regions = 1}, we associate the Physical Line 1, defined in the Step 1, to the \TiberCAD{} \textbf{Region} \textit{bulk}.

\subsection* {2 - Definition of Simulation}

Now we define the \textbf{Simulation} \textit{driftdiffusion\_1}: it belongs to  the class \textbf{driftdiffusion}
\begin{verbatim}
Models
{
model driftdiffusion
{
options
{
simulation_name = driftdiffusion_1
physical_regions = all
}
\end{verbatim}

The \TiberCAD{} \textbf{simulation} \textit{driftdiffusion\_1} , belonging to the  \textbf{model} \textbf{driftdiffusion}, will be applied to the whole device structure (\textit{physical\_regions = all}).


\subsection* {3 - Definition of Boundary Conditions}
The anode and cathode contacts of our 1D Si sample are defined as \textbf{Boundary conditions regions} (\textit{BC\_Region anode , BC\_Region cathode}) in the following way:

\begin{verbatim}
BC_Region anode
{
BC_reg_numb = 1
type = ohmic
voltage = @Vb
}

BC_Region cathode
{
BC_reg_numb = 2
type = ohmic
voltage = 0.0
}
\end{verbatim}
Both contacts are defined as \textit{ohmic}, cathod is assigned a fixed \textit{voltage = 0.0}, while anode voltage is given by the value of the variable \textit{Vb (voltage = @Vb)}.

Through the keyword  \textbf{BC\_reg\_numb} , one or more of the \textbf{(n-1)D} physical regions (Physical Points in 1D, Physical Lines in 2D, Physical Surfaces in 3D) previously defined in the GMSH mesh can be associated to the present  \TiberCAD{} \textbf{BC Region}.
With \textit{BC\_reg\_numb = 1}, we associate the Physical Point 1, defined in the \textbf{Step 1}, to the \TiberCAD{} \textbf{BC Region} \textit{anode}; with \textit{BC\_reg\_numb = 2}, we associate the Physical Point 2, defined in the \textbf{Step 1}, to the \TiberCAD{} \textbf{BC Region} \textit{cathode}.

Alternatively, one can make  use  of  the  \textbf{physical names} associated to the physical regions in the  meshing  tool.
In  this  case, we  simply associate the  \textbf{(n-1)D} physical region,  respectively  ``anode'' and  ``cathode'',  by  means  of  the  \TiberCAD{} \textbf{BC Region name}:

 \begin{verbatim}
BC_Region anode
{

type = ohmic
voltage = @Vb
}

BC_Region cathode
{

type = ohmic
voltage = 0.0
}
\end{verbatim}

Note that, in  this  case, the \TiberCAD{} \textbf{BC Region name} needs to  be  identical to  one of the  \textbf{physical names} defined during the  modeling of the  device  with   GMSH.



\subsection* {4 - Definition of Simulation parameters}

The variable \textit{Vb }is specified in the \textit{sweep} block, in the \textbf{Solver} section
\begin{verbatim}
 sweep
{
simulation = driftdiffusion_1
variable = Vb
start = 0.0
stop = 1
steps = 10
}
\end{verbatim}

In this way, the simulation \textit{driftdiffusion\_1} is performed for 10 ( \textit{steps = 10}) values of the anode voltage (\textit{variable = Vb}), between 0 and 1.

\subsection* {5 - Definition of Execution parameters}
In the \textbf{Simulation} section , we decide \textit{which} simulations to perform and in which \textit{order}; we set
\textit{solve = sweep}, to execute the \textit{sweep} which run \textit{driftdiffusion\_1} \textbf{simulation} for the specified loop.
\begin{verbatim}
  $Simulation
{
#searchpath = .
meshfile = bulk.msh
dimension = 1
temperature = 300
solve = sweep
resultpath = output
output_format = grace
plot = (Ec, Ev, ContactCurrents)
}
\end{verbatim}
Output files with conduction and valence band profiles (\textit{plot = Ec,Ev.}.) and all the calculated values of the current at the contacts (\textit{ContactCurrents}) (the IV characteristic) are generated.



\section*{Step 4: Run TiberCAD{}}

Now we can run TiberCAD{}:

\texttt{tibercad  bulk.tib}


\vspace{0.5cm}
The generated \textbf{Output} files are:



\textbf{driftdiffusion\_materials.dat}: \textit{material} (mesh) regions, in this case just region 1

\textbf{driftdiffusion\_nodal.dat}: \textit{nodal} quantities (here conduction and valence band)

\textbf{sweep\_driftdiffusion\_Vb.dat} : \textit{integrated current} at the two contacts for each sweep step.














